% ============================================================
%  Plano Cartesiano — Apostila Premium | PratiqueJá
%  Engine: XeLaTeX
% ============================================================
\documentclass[12pt]{article}
\usepackage{../../pratiqueja}
\usepackage{tikz}

% \hlm — destaque de resultado em modo-math (cor sem bold)
\newcommand{\hlm}[1]{{\color{darkblue}#1}}

\tikzset{
  axe/.style={->, line width=1pt, gray!60!black},
  gridln/.style={very thin, gray!25},
  pt/.style={circle, fill=darkblue, inner sep=1.6pt},
  ptlbl/.style={font=\footnotesize\bfseries, darkblue},
  quad/.style={font=\large\bfseries, gray!45},
}

\setsubject{Plano Cartesiano}

\begin{document}
\pagenumbering{arabic}
\setstretch{1.2}
\color{bodytext}

\heroheader{Plano Cartesiano}%
           {Eixos, Quadrantes, Pontos e Simetria}%
           {Matemática · Básico}

\setlength{\columnsep}{18pt}
\setlength{\columnseprule}{0.5pt}
\def\columnseprulecolor{\color{exborder}}

\begin{multicols}{2}

%─────────────────────────────────────────────────────────────
\TW{Def}{badgepurple}{Eixos e Origem}{A estrutura do plano e os quatro quadrantes}

\regra{O \textbf{plano cartesiano} é formado por duas retas numéricas
perpendiculares que se cruzam na \textbf{origem}:\\[3pt]
\textbf{Eixo $x$} (horizontal) — das \emph{abscissas}\\
\textbf{Eixo $y$} (vertical) — das \emph{ordenadas}\\
\textbf{Origem} $O(0,0)$ — interseção dos eixos}

\begin{center}\vspace{1pt}
\begin{tikzpicture}[scale=0.52]
% grid
\draw[gridln] (-3.4,-3.4) grid (4.4,3.4);
% eixos
\draw[axe] (-3.6,0) -- (4.7,0) node[right, font=\small\bfseries, black]{$x$};
\draw[axe] (0,-3.6) -- (0,3.7) node[above, font=\small\bfseries, black]{$y$};
\node[below left, font=\footnotesize] at (0,0) {$O$};
% rótulos de quadrante
\node[quad] at (2.4,2.4) {I};
\node[quad] at (-2.4,2.4) {II};
\node[quad] at (-2.4,-2.4) {III};
\node[quad] at (2.4,-2.4) {IV};
% pontos
\node[pt] at (2,2) {};   \node[ptlbl, above right] at (2,2) {A(2,2)};
\node[pt] at (-1,1) {};  \node[ptlbl, above left] at (-1,1) {B(-1,1)};
\node[pt] at (-2,-1) {}; \node[ptlbl, below left] at (-2,-1) {C(-2,-1)};
\node[pt] at (3,-2) {};  \node[ptlbl, below right] at (3,-2) {D(3,-2)};
\end{tikzpicture}
\end{center}

\SH{Os quatro quadrantes (anti-horário)}
\begin{center}\vspace{1pt}
\setlength{\tabcolsep}{6pt}\renewcommand{\arraystretch}{1.25}\small
\begin{tabular}{c c c c}
\rowcolor{rulebg}\textbf{Quad.} & \textbf{$x$} & \textbf{$y$} & \textbf{Ex.}\\
\textbf{I}   & $+$ & $+$ & A(2,2)\\
\textbf{II}  & $-$ & $+$ & B(-1,1)\\
\textbf{III} & $-$ & $-$ & C(-2,-1)\\
\textbf{IV}  & $+$ & $-$ & D(3,-2)\\
\end{tabular}
\end{center}

\nota{Pontos \textbf{sobre os eixos} não pertencem a quadrante algum:
$(3,0)$ está no eixo $x$ e $(0,-2)$ no eixo $y$.}

%─────────────────────────────────────────────────────────────
\TW{Loc}{darkblue}{Localização de Pontos}{Par ordenado $(x,y)$ — abscissa e ordenada}

\regra{Cada ponto é um \textbf{par ordenado} $(x,y)$:\\[3pt]
\textbf{Abscissa} $(x)$ — distância horizontal (direita $+$, esq.\ $-$)\\
\textbf{Ordenada} $(y)$ — distância vertical (cima $+$, baixo $-$)\\[3pt]
A ordem importa: $(2,3) \neq (3,2)$.}

\exemp{%
  \textbf{Marcar $P(3,-1)$:}\\
  1. da origem, \hlm{3 à direita} $(x=3)$\\
  2. depois \hlm{1 para baixo} $(y=-1)$\\
  3. marque $P$ nessa posição%
}

\exemp{%
  \textbf{Em que quadrante está $Q(-4,2)$?}\\
  $x=-4$ (neg.) e $y=2$ (pos.) → \hlm{Quadrante II}%
}

%─────────────────────────────────────────────────────────────
\TW{Dist}{badgegreen}{Distância à Origem}{Teorema de Pitágoras no plano}

\regra{Os catetos são as projeções $|x|$ e $|y|$:}
\formula{$d(P,O) = \sqrt{x^2 + y^2}$}

\exemp{%
  \textbf{$A(3,4)$:}\\
  $d = \sqrt{3^2+4^2} = \sqrt{9+16} = \sqrt{25} = \hlm{5}$\\[3pt]
  \textbf{$B(-1,-1)$:}\\
  $d = \sqrt{(-1)^2+(-1)^2} = \sqrt{2} \approx \hlm{1{,}41}$\\[3pt]
  \textbf{$C(5,0)$ (sobre o eixo $x$):}\\
  $d = \sqrt{5^2+0^2} = \hlm{5} = |x|$%
}

%─────────────────────────────────────────────────────────────
\TW{Sim}{badgegreen}{Pares Simétricos}{Reflexão nos eixos e na origem}

\regra{Dado $P(a,b)$, seus simétricos são:\\[3pt]
$P'_x = (a,\,-b)$ — reflexão no eixo $x$\\
$P'_y = (-a,\,b)$ — reflexão no eixo $y$\\
$P'_O = (-a,\,-b)$ — reflexão na origem}

\exemp{%
  \textbf{$P(3,2)$:}\\
  eixo $x$: $\hlm{(3,-2)}$\\
  eixo $y$: $\hlm{(-3,2)}$\\
  origem: $\hlm{(-3,-2)}$%
}

\exemp{%
  \textbf{$Q(-4,1)$:}\\
  eixo $x$: $(-4,-1)$ — Quadrante III\\
  eixo $y$: $(4,1)$ — Quadrante I\\
  origem: $(4,-1)$ — Quadrante IV%
}

\nota{O simétrico ao eixo $x$ mantém a abscissa e inverte a ordenada.
O simétrico à origem inverte \emph{os dois} sinais — equivale a duas
reflexões seguidas.}

\end{multicols}

\end{document}
